\chapter{kbctl Reference}
\label{ch:kbctl}
\cmd{kbctl} is the programmatic workhorse. Every control operation is a typed protobuf gRPC request over
\file{/run/kb/kba.sock}. The reference text below is \textbf{generated from the installed binary's
\texttt{--help}} by \file{latex/field-manual/gen-help.sh}, so it cannot drift from the code.

\newcommand{\helpblock}[2]{\subsection{#1}\lstinputlisting[style=console]{generated/kbctl-#2.txt}}

\section{Command map}
\begin{figure}[ht]
\centering
\begin{tikzpicture}[x=1cm,y=1cm,scale=0.95,transform shape]
  \node[dark,minimum width=2.4cm,minimum height=1.1cm] (root) at (0,-5.4) {kbctl};
  \node[soft,minimum width=3.6cm] (g0) at (4,-0.0) {\texttt{up | down | status}};
  \node[anchor=west,font=\bfseries\sffamily\footnotesize] at (6.1,-0.0) {whole-stack shortcuts};
  \node[dark,minimum width=3.6cm] (g1) at (4,-1.1) {dashboard};
  \node[anchor=west,font=\ttfamily\bfseries\footnotesize] at (6.1,-1.1) {up | down | status | url | open};
  \node[dark,minimum width=3.6cm] (g2) at (4,-2.2) {ssh};
  \node[anchor=west,font=\ttfamily\bfseries\footnotesize] at (6.1,-2.2) {up | down | status | connect | add-key};
  \node[dark,minimum width=3.6cm] (g3) at (4,-3.3000000000000003) {tui};
  \node[anchor=west,font=\ttfamily\bfseries\footnotesize] at (6.1,-3.3000000000000003) {runs kbopd on this terminal};
  \node[dark,minimum width=3.6cm] (g4) at (4,-4.4) {process};
  \node[anchor=west,font=\ttfamily\bfseries\footnotesize] at (6.1,-4.4) {isolate};
  \node[dark,minimum width=3.6cm] (g5) at (4,-5.5) {zone};
  \node[anchor=west,font=\ttfamily\bfseries\footnotesize] at (6.1,-5.5) {override};
  \node[dark,minimum width=3.6cm] (g6) at (4,-6.6000000000000005) {policy};
  \node[anchor=west,font=\ttfamily\bfseries\footnotesize] at (6.1,-6.6000000000000005) {reload};
  \node[dark,minimum width=3.6cm] (g7) at (4,-7.700000000000001) {workload};
  \node[anchor=west,font=\ttfamily\bfseries\footnotesize] at (6.1,-7.700000000000001) {keygen | sign | reload};
  \node[dark,minimum width=3.6cm] (g8) at (4,-8.8) {audit};
  \node[anchor=west,font=\ttfamily\bfseries\footnotesize] at (6.1,-8.8) {verify | export};
  \node[dark,minimum width=3.6cm] (g9) at (4,-9.9) {stats};
  \node[anchor=west,font=\ttfamily\bfseries\footnotesize] at (6.1,-9.9) {(no subcommands)};
  \foreach \i in {0,...,9} \draw[line width=0.8pt] (root.east) -- ++(0.7,0) |- (g\i.west);
\end{tikzpicture}
\caption{Everything under \texttt{kbctl}. The first row are top-level shortcuts that act on the whole stack.}
\end{figure}

\section{Global flag}
\begin{lstlisting}[style=console]
--socket string   path to kbd's gRPC Unix domain socket (default "/run/kb/kba.sock")
\end{lstlisting}
Only commands that talk to \cmd{kbd} use it (\texttt{policy}, \texttt{workload reload}, \texttt{zone},
\texttt{process}, \texttt{audit}, \texttt{stats}). The service-control commands (\texttt{up}, \texttt{down},
\texttt{status}, \texttt{dashboard}, \texttt{ssh}) go through \cmd{systemctl} and work even when \cmd{kbd} is down.

\section{Exit codes}
\begin{center}\small
\begin{tabularx}{\textwidth}{@{}L{1.6cm}Y@{}}
\toprule
\textbf{Code} & \textbf{Meaning}\\
\midrule
0 & Success. For \texttt{dashboard status} and \texttt{ssh status}: the unit is active \emph{and} the port answers.\\
1 & Error (printed as \texttt{error: ...}): bad arguments, RPC failure, \texttt{systemctl} failure.\\
3 & \texttt{dashboard status} / \texttt{ssh status} only: unit not active or port not listening. Use it in scripts.\\
\bottomrule
\end{tabularx}
\end{center}

\section{Stack control}
\helpblock{kbctl up}{up}
\helpblock{kbctl down}{down}
\helpblock{kbctl status}{status}

\section{Web dashboard}
\helpblock{kbctl dashboard}{dashboard}
\helpblock{kbctl dashboard open}{dashboard-open}
\helpblock{kbctl dashboard status}{dashboard-status}
\helpblock{kbctl dashboard url}{dashboard-url}

\section{SSH console and TUI}
\helpblock{kbctl ssh}{ssh}
\helpblock{kbctl ssh connect}{ssh-connect}
\helpblock{kbctl ssh add-key}{ssh-add-key}
\helpblock{kbctl tui}{tui}

\section{Containment and classification}
\helpblock{kbctl process isolate}{process-isolate}
\textbf{Levels} are \texttt{CGROUP}, \texttt{SECCOMP}, \texttt{NAMESPACE} (default) and \texttt{TERMINATE}, in
increasing severity. The command prints \texttt{pid=N contained at level L}. The reason is written to the audit log.
\helpblock{kbctl zone override}{zone-override}
\begin{kbnote}[OVERRIDE IS COSMETIC, NOT ENFORCEMENT]
\texttt{zone override} relabels how \cmd{kbd} \emph{displays and classifies} a process. It does not touch
kernel or enforcement state, and the next real zone transition from the sensor overwrites it. To actually
restrict a process, use \texttt{process isolate}.
\end{kbnote}

\section{Policy, workloads, audit, stats}
\helpblock{kbctl policy reload}{policy-reload}
\helpblock{kbctl workload keygen}{workload-keygen}
\helpblock{kbctl workload sign}{workload-sign}
\helpblock{kbctl workload reload}{workload-reload}
\helpblock{kbctl audit verify}{audit-verify}
\helpblock{kbctl audit export}{audit-export}
\helpblock{kbctl stats}{stats}

\section{The daemons' own flags}
\subsection{kbd}
\lstinputlisting[style=console]{generated/kbd.txt}
\subsection{kb-checker (kept down; for reference)}
\lstinputlisting[style=console]{generated/kb-checker-monitor.txt}
